\documentclass[11pt,a4paper]{report}

% ---------------------------------------------------------------------------
% Packages (conservative, widely available in TeX Live)
% ---------------------------------------------------------------------------
\usepackage[T1]{fontenc}
\usepackage[utf8]{inputenc}
\usepackage{amsmath}
\usepackage{amssymb}
\usepackage[a4paper,margin=2.4cm]{geometry}
\usepackage{xcolor}
\usepackage{listings}
\usepackage{booktabs}
\usepackage{longtable}
\usepackage{array}
\usepackage{enumitem}
\usepackage{fancyhdr}
\usepackage{parskip}
\usepackage[hidelinks]{hyperref}

% ---------------------------------------------------------------------------
% Colours & styles
% ---------------------------------------------------------------------------
\definecolor{accent}{HTML}{5B21B6}
\definecolor{accent2}{HTML}{7C3AED}
\definecolor{codebg}{HTML}{F5F5FA}
\definecolor{codeborder}{HTML}{D0D0E0}
\definecolor{kw}{HTML}{7C3AED}
\definecolor{str}{HTML}{16A34A}
\definecolor{com}{HTML}{6B7280}
\definecolor{zgreen}{HTML}{16A34A}
\definecolor{zamber}{HTML}{D97706}
\definecolor{zred}{HTML}{DC2626}
\definecolor{muted}{HTML}{555577}

\hypersetup{colorlinks=true, linkcolor=accent, urlcolor=accent2}

\lstdefinestyle{py}{
  language=Python,
  backgroundcolor=\color{codebg},
  basicstyle=\ttfamily\small,
  keywordstyle=\color{kw}\bfseries,
  stringstyle=\color{str},
  commentstyle=\color{com}\itshape,
  frame=single, rulecolor=\color{codeborder},
  showstringspaces=false,
  breaklines=true, breakatwhitespace=true,
  columns=fullflexible, keepspaces=true,
  numbers=none, xleftmargin=6pt, xrightmargin=4pt,
  aboveskip=8pt, belowskip=6pt,
}
\lstdefinestyle{plain}{
  backgroundcolor=\color{codebg},
  basicstyle=\ttfamily\small,
  frame=single, rulecolor=\color{codeborder},
  breaklines=true, columns=fullflexible, keepspaces=true,
  xleftmargin=6pt, xrightmargin=4pt, aboveskip=8pt, belowskip=6pt,
}
\lstset{style=py}
\newcommand{\code}[1]{\texttt{#1}}

% Headers / footers
\pagestyle{fancy}
\fancyhf{}
\lhead{\small\textcolor{muted}{Wafer Map Tool}}
\rhead{\small\textcolor{muted}{Architecture \& Design}}
\rfoot{\small\thepage}
\lfoot{\small\textcolor{muted}{\leftmark}}
\renewcommand{\headrulewidth}{0.3pt}

% ---------------------------------------------------------------------------
\begin{document}

% ===== Title page ==========================================================
\begin{titlepage}
\centering
\vspace*{3cm}
{\Huge\bfseries\textcolor{accent}{Wafer Map Tool}}\\[6pt]
{\LARGE Architecture, Pipeline, GUI \& Automation}\\[4pt]
{\large Technical Reference}\\[2cm]
\begin{tabular}{rl}
\textbf{Component} & Wafer Map Tool (\code{wafer\_tool} package)\\[4pt]
\textbf{Scope} & EFF ingestion, wafer-map rendering, PDF/PPTX/CSV\\
 & reporting, interactive GUI, offline automation\\[4pt]
\textbf{Platform} & Windows (PyQt6); headless runner is cross-platform\\[4pt]
\textbf{Rendering} & Matplotlib (Agg) + SciPy interpolation\\
\end{tabular}
\vfill
{\large\textcolor{muted}{Generated \today}}
\end{titlepage}

\tableofcontents
\clearpage

% ===========================================================================
\chapter{Introduction}

The \textbf{Wafer Map Tool} turns per-die parametric measurements into
colour-coded wafer maps and packages them into a multi-page \textbf{PDF},
a \textbf{PowerPoint} deck, and a summary \textbf{CSV}. It accepts two kinds of
input:

\begin{itemize}[leftmargin=1.4em]
  \item \textbf{Raw \code{.eff} extraction files} — the native semicolon-delimited
        export from the eSquare / EBS-XE test flow, containing many measurement
        parameters. The tool scans the header, lets the user pick which
        parameters to map, and produces one output folder per parameter.
  \item \textbf{Converted \code{.txt}/\code{.csv} files} — a single measurement
        column already extracted to a four-column
        \code{lot wafer\,/\,x\,/\,y\,/\,value} layout.
\end{itemize}

The application has three faces sharing one rendering core:

\begin{enumerate}[leftmargin=1.6em]
  \item An interactive \textbf{PyQt6 GUI} with live wafer-thumbnail preview,
        threshold scatter and histogram tabs, and a per-parameter preview pager.
  \item A \textbf{headless runner} that replays a saved job with no GUI, suitable
        for scheduled / offline execution.
  \item A \textbf{Windows Task Scheduler} integration that runs saved jobs
        daily / weekly / monthly and reports their status.
\end{enumerate}

\section{Reading this document}
Chapter~\ref{ch:arch} gives the module map and layering rules. Chapters
\ref{ch:data}--\ref{ch:exporters} follow the data through the pipeline, from file
formats to the rendered report. Chapter~\ref{ch:gui} covers the GUI and
Chapter~\ref{ch:auto} the automation. Chapters
\ref{ch:concurrency}--\ref{ch:modularity} discuss the concurrency model,
performance and the modularity assessment.

% ===========================================================================
\chapter{High-Level Architecture}\label{ch:arch}

\section{Layered design}
The package is organised as a set of layers with a strict downward dependency
rule: higher layers may import lower ones, never the reverse. The GUI and the
headless runner are two \emph{drivers} that sit on top of the same
GUI-free core.

\begin{lstlisting}[style=plain]
+---------------------------------------------------------------+
|  DRIVERS                                                      |
|  ui/ (PyQt6 widgets, main window)   automation/runner (CLI)   |
+------------------------+--------------------------------------+
|  ORCHESTRATION         |                                      |
|  report.generate_report        eff_pipeline.process_eff       |
|  eff_worker (Qt)  worker (Qt)   automation/job, scheduler      |
+---------------------------------------------------------------+
|  EXPORTERS                                                    |
|  exporters/pdf_exporter  pptx_exporter  csv_exporter          |
|  exporters/report_charts (scatter / histogram pages)          |
+---------------------------------------------------------------+
|  CORE RENDERING & DATA                                        |
|  plotting  geometry  statistics  data_loader  eff_loader      |
|  config (PlotConfig, ColorScheme)   logger   exceptions       |
+---------------------------------------------------------------+
             (compiled C parser: eff_parser is external)
\end{lstlisting}

\section{Module responsibilities}
\begin{longtable}{@{}p{4.3cm}p{10.2cm}@{}}
\toprule
\textbf{Module} & \textbf{Responsibility}\\
\midrule
\endhead
\code{config.py} & \code{PlotConfig} (thresholds, log, mirror, rotation) and
  \code{ColorScheme}; tuning constants (\code{GRID\_RESOLUTION},
  \code{MAX\_WAFERS\_PER\_PAGE}, \code{PAGE\_PNG\_DPI}).\\
\code{geometry.py} & Spatial transforms (mirror / rotation) and
  \code{build\_wafer\_grid} — the SciPy \code{griddata} interpolation onto a
  square mesh.\\
\code{plotting.py} & \code{draw\_wafer\_ax} — renders one wafer contour map onto a
  Matplotlib axes.\\
\code{statistics.py} & 8-condition classification counts and the summary bar
  chart page.\\
\code{data\_loader.py} & Reads converted \code{.txt}/\code{.csv} into a typed
  DataFrame (\code{lot, wafer, x, y, value}).\\
\code{eff\_loader.py} & Self-contained \code{.eff} reader: header scan, coordinate
  / measurement-column resolution, spec-limit capture, streaming per-parameter
  extraction, and an in-memory single-parameter reader for previews.\\
\code{report.py} & End-to-end pipeline orchestrator: load $\to$ statistics $\to$
  per-wafer PNGs $\to$ PDF $\to$ CSV $\to$ PPTX.\\
\code{eff\_pipeline.py} & GUI-free orchestration for the multi-parameter EFF
  workflow (\code{process\_eff}); shared by the GUI worker and the headless
  runner.\\
\code{eff\_worker.py} & Thin Qt \code{QThread} adapters: \code{EffReportWorker}
  and \code{EffPreviewWorker}.\\
\code{worker.py} & \code{ReportWorker} \code{QThread} for the single-file
  (txt/csv) workflow.\\
\code{exporters/} & \code{pdf\_exporter}, \code{pptx\_exporter},
  \code{csv\_exporter}, and \code{report\_charts} (overview scatter/histogram).\\
\code{ui/} & \code{main\_window}, \code{histogram\_widget},
  \code{eff\_param\_dialog}, \code{orientation\_widget}, \code{threshold\_viz\_widget}.\\
\code{automation/} & \code{job} (the \code{.wtjob} format), \code{runner}
  (headless replay), \code{scheduler} (Task Scheduler), \code{scheduler\_dialog}.\\
\bottomrule
\end{longtable}

\section{Key design principle}
The most important structural decision is that \textbf{all report generation is
GUI-free}. \code{report.generate\_report} and \code{eff\_pipeline.process\_eff}
depend only on Matplotlib/pandas/SciPy, never on PyQt. This is what makes the
same logic runnable both interactively (behind a \code{QThread}) and headless
(from the scheduler), guaranteeing identical output.

% ===========================================================================
\chapter{Data Model \& File Formats}\label{ch:data}

\section{The EFF extraction file}
An \code{.eff} file is a \emph{semicolon-delimited text} export. After a block of
\code{<<Section>>} metadata it carries a set of \code{<...>} header rows and then
one row per die:

\begin{lstlisting}[style=plain]
<<EFF:1.00>>;Headers=24;Rows=1527;Columns=109;...
<DataType>;Text;Int;Int;Int;...;Double;Double;...
<+ParameterName>;Lot;Wafer;X;Y;...;RON_VD__10_75;...
<SourceTable>;;;MeasPartPos;MeasPartPos;...;MeasPartVals;...
<LIMIT:SPEC:LOWER_VALUE>;;;;;...;0.0;...
<LIMIT:SPEC:UPPER_VALUE>;;;;;...;31.0;...
1E351778;01;8;5;...;-113.25;...        <- data rows
\end{lstlisting}

\noindent Resolution rules (implemented in \code{eff\_loader.py}, mirroring the
standalone \code{eff\_converter}):
\begin{itemize}[leftmargin=1.4em]
  \item \textbf{Coordinates} — the columns named \code{Lot} (or \code{LotNumber}),
        \code{Wafer}, \code{X}, \code{Y} are located case-insensitively and are
        always loaded.
  \item \textbf{Measurement parameters} — a column is selectable when its
        \code{<SourceTable>} is \code{MeasPartVals}, or (fallback) it is a
        floating type carrying spec limits.
  \item \textbf{Spec limits} — \code{<LIMIT:SPEC:LOWER\_VALUE>} and
        \code{UPPER\_VALUE} are captured per parameter and used to pre-fill the
        GUI thresholds.
\end{itemize}

The row count (\code{Rows=}) is read from the \code{<<EFF>>} line for a progress
estimate; only the header is read during a \emph{scan}, so scanning a multi-GB
file is fast.

\section{The converted TXT layout}
The single-parameter text file is four tab-separated columns:
\begin{lstlisting}[style=plain]
1E351778 01<TAB>8<TAB>5<TAB>-113.25
   |          |    |    |
   lot wafer  x    y    value
\end{lstlisting}
\code{data\_loader.read\_wafer\_data} parses this with pandas' fast C engine,
forcing the identifier column to text (so leading-zero wafer ids like ``01''
survive) and letting numeric columns parse natively.

\section{PlotConfig \& ColorScheme}
\code{PlotConfig} is the single typed bundle of user plotting choices passed
through the whole pipeline:

\begin{lstlisting}
@dataclass
class PlotConfig:
    t_low: float          # Limit 1
    t_high: float         # Limit 2  (auto-ordered t_low <= t_high)
    use_log: bool = False
    high_is_green: bool = False
    mirror_x: bool = False
    mirror_y: bool = False
    rot_deg: int = 0      # one of 0/90/180/270
\end{lstlisting}

\code{ColorScheme.from\_mode(high\_is\_green)} yields the three zone colours
(\textcolor{zgreen}{low}/\textcolor{zamber}{mid}/\textcolor{zred}{high} or the
reverse). Zones: value $<$ Limit\,1, Limit\,1 $\le$ value $\le$ Limit\,2, and
value $>$ Limit\,2.

\section{The job file (.wtjob)}
A \code{.wtjob} is JSON capturing everything ``Generate Report'' needs, so a run
can be replayed offline. Beyond the \code{PlotConfig} fields it stores:

\begin{longtable}{@{}p{3.6cm}p{10.9cm}@{}}
\toprule
\textbf{Field} & \textbf{Meaning}\\
\midrule
\endhead
\code{input\_file} & Fixed input path (txt/csv or \code{.eff}).\\
\code{input\_dir} + \code{input\_pattern} & Folder mode: pick the newest file
  matching the pattern at run time (e.g.\ \code{*.eff}) — so a scheduled job
  always processes the latest drop.\\
\code{out\_dir} & Output directory.\\
\code{eff\_params} & For raw \code{.eff}: the parameter names to map (one folder
  each). Empty means ``all parameters''.\\
\code{eff\_filters} & Optional per-parameter value filter: parameter name
  $\to [\min, \max]$. Name-keyed so a folder-mode job re-maps it onto the newest
  file. Empty means no filter.\\
\bottomrule
\end{longtable}

% ===========================================================================
\chapter{The Report Pipeline}\label{ch:pipeline}

\section{\code{generate\_report}}
The single-file pipeline is \code{report.generate\_report(filepath, config,
out\_dir, on\_progress, on\_wafer\_ready, render\_individual\_pngs, jobs)}. Its
stages:

\begin{enumerate}[leftmargin=1.6em]
  \item \textbf{Load} the data (\code{read\_wafer\_data}).
  \item \textbf{Statistics} — \code{calculate\_8\_condition\_statistics} classifies
        every die into one of eight conditions.
  \item \textbf{Per-wafer PNGs} (optional) — one square archive image per wafer,
        rendered in parallel; feeds the live GUI preview and lands in the output.
  \item \textbf{PDF} — \code{generate\_pdf} builds the multi-page report.
  \item \textbf{CSV} — \code{export\_color\_summary\_csv}.
  \item \textbf{PPTX} — \code{create\_powerpoint\_report} assembles page snapshots.
\end{enumerate}

Two callbacks decouple the pipeline from any UI: \code{on\_progress(current,
total, message)} and \code{on\_wafer\_ready(lot, wafer, png\_path)}. A
\code{RuntimeError} raised from \code{on\_progress} is the cancellation channel.

\section{Progress \& cancellation contract}
\code{total\_steps = n\_wafers + 3} (three export steps). The wafer-render phase
reports one step per wafer; the last three are PDF / CSV / PPTX. The GUI worker
translates these into two progress bars and cancels by raising
\code{RuntimeError("Cancelled by user.")} from its progress slot.

% ===========================================================================
\chapter{EFF Ingestion Pipeline}\label{ch:eff}

The EFF workflow turns \emph{one} file with \emph{many} parameters into one
report per selected parameter. It is deliberately split so the interactive and
automated paths run identical logic.

\section{Scan}
\code{eff\_loader.scan\_eff(path)} reads only the header and returns an
\code{EffScan}: declared row/column counts, the resolved coordinate indices, and
a list of \code{EffParameter} (index, name, dtype, \code{has\_limits},
\code{limit\_lower}, \code{limit\_upper}). It stops at the first data row, so it
is fast on very large files.

\section{Extraction}
\code{extract\_parameters(path, indices, out\_base\_dir, ...)} streams the whole
file \emph{once} and writes one four-column text file per selected parameter into
its own folder:
\begin{lstlisting}[style=plain]
<out_dir>/<parameter>/<parameter>.txt
\end{lstlisting}
Invalid values (empty, non-numeric, or $|v|\ge 10^{10}$ sentinels) are skipped.
The per-parameter folder subsequently receives that parameter's PNGs, PDF, PPTX
and CSV.

\section{Orchestration: \code{process\_eff}}
\code{eff\_pipeline.process\_eff} is the GUI-free core:
\begin{enumerate}[leftmargin=1.6em]
  \item Extract all selected parameters in one streaming pass
        (\textasciitilde30\% of the progress bar).
  \item For each parameter that yielded data, call \code{generate\_report} into
        that parameter's folder.
  \item Flatten the per-wafer PNGs up into the folder so images + PDF + PPTX + CSV
        sit together.
\end{enumerate}
Progress is reported as \code{(overall\_pct, wafer\_pct, message)}; cancellation
is an \code{EffCancelled} exception (converted to a \code{RuntimeError} inside the
\code{generate\_report} call so the render aborts cleanly).

\section{Per-parameter value filter}
Extraction accepts an optional \code{filters} mapping
(\code{\{column\_index: (min, max)\}}); a die whose value falls outside a
parameter's range is excluded from that parameter's output --- alongside the
always-on $|v|\ge 10^{10}$ sentinel. The parameter picker defaults every
row to a universal $\pm1000$ range (matching the standalone eff-converter),
editable per parameter, and the whole filter can be switched off. The same range
is applied to the preview reader, so the preview matches the exported map die for
die. The filter controls \emph{inclusion}; it is independent of the colour
thresholds (Limit\,1 / Limit\,2), which only control zone colouring. For
scheduled jobs the ranges are persisted name-keyed in \code{eff\_filters} and
re-mapped onto whatever parameters the newest file contains.

\section{In-memory preview reader}
\code{extract\_single\_param\_df} reads exactly one parameter into a DataFrame
without touching disk. The GUI uses it (on a background thread) to draw the
scatter/histogram preview for the currently-selected parameter, honouring the
value filter above.

% ===========================================================================
\chapter{Rendering Internals}\label{ch:render}

\section{Interpolation (\code{geometry.build\_wafer\_grid})}
Each wafer's scattered $(x,y,\text{value})$ dies are interpolated onto a
\code{GRID\_RESOLUTION}$^2$ (default $100\times100$) square mesh with SciPy
\code{griddata}, clamped to the data range, then masked to the wafer disc. This
is the CPU-heavy step.

\section{Drawing (\code{plotting.draw\_wafer\_ax})}
Applies the spatial transforms, interpolates, maps values through the
log-or-linear threshold levels, and paints three filled contour bands with the
\code{ColorScheme} colours plus the wafer-edge circle.

\section{Parallel rendering}
Per-wafer PNGs and PDF grid pages are independent, so both render across a
\code{ProcessPoolExecutor}. On Windows each worker re-imports Matplotlib
(\textasciitilde7\,s fixed spawn cost), so pools are only used when they pay off:

\begin{longtable}{@{}p{6.5cm}p{8.0cm}@{}}
\toprule
\textbf{Pass} & \textbf{Parallel threshold}\\
\midrule
\endhead
Per-wafer PNGs (\code{report.\_render\_wafer\_pngs}) & $\ge 200$ wafers
  (\code{\_PNG\_PARALLEL\_MIN}).\\
PDF grid pages (\code{pdf\_exporter.\_render\_all\_grid\_pages}) & $\ge 8$ pages
  $\approx 200$ wafers (\code{\_PDF\_PARALLEL\_MIN\_PAGES}).\\
\bottomrule
\end{longtable}

\noindent Measured crossover is \textasciitilde160 wafers (100 wafers: parallel
is $0.57\times$, i.e.\ slower; 500 wafers: $1.94\times$ faster). With both passes
parallel, a 500-wafer report improved from \textasciitilde55\,s to
\textasciitilde25\,s ($\sim2.2\times$).

\section{Adaptive per-wafer DPI}
The individual per-wafer archive PNGs use an adaptive output resolution keyed to
how many threshold zones the wafer's die values span (\code{report.\_zone\_count},
computed from the raw values --- no interpolation needed):
\begin{itemize}[leftmargin=1.4em]
  \item a wafer entirely within one zone renders as a single flat colour and is
        saved at \code{WAFER\_PNG\_DPI\_UNIFORM} (70\,dpi);
  \item a wafer spanning more than one zone has colour boundaries worth resolving
        and is saved at \code{WAFER\_PNG\_DPI\_MULTIZONE} (150\,dpi).
\end{itemize}
Measured: a uniform wafer $\to 560\times560$\,px ($\sim$5.6\,KB); a three-zone
wafer $\to 1200\times1200$\,px ($\sim$14\,KB). The rule keys off the zone count,
not the specific colour, so an all-yellow or all-red wafer is also treated as
uniform. The PDF grid cells are vector (resolution-independent) and the PPTX page
snapshots mix 25 wafers, so adaptive DPI applies only to the standalone
per-wafer images.

\section{Overview charts (\code{exporters/report\_charts.py})}
Headless, Matplotlib-only renderers for the report's front-page \textbf{Threshold
Scatter} and \textbf{Value Histogram}, colour-matched to the wafer maps:
\begin{itemize}[leftmargin=1.4em]
  \item Scatter plots value against wafer (jittered into a per-wafer strip),
        downsampling above 200{,}000 points.
  \item Each threshold zone gets a distinct \textbf{shape}
        (\textcolor{zred}{$\blacklozenge$ diamond} below Limit\,1,
        \textcolor{zamber}{$\blacksquare$ square} mid,
        \textcolor{zgreen}{$\bullet$ circle} above Limit\,2).
  \item Marker size/alpha are \textbf{adaptive}: big and bold for sparse
        parameters, small and light for dense ones (100k+ dies) so they do not
        smear into solid bands ($\sim1$k pts $\to s\approx20$; $\sim200$k pts
        $\to s\approx3$).
\end{itemize}

% ===========================================================================
\chapter{Exporters}\label{ch:exporters}

\section{PDF structure}
\code{pdf\_exporter.generate\_pdf} builds a head section in the main process and
merges per-lot grid pages (rendered in worker processes) with \code{pypdf},
attaching bookmarks. Page order:

\begin{enumerate}[leftmargin=1.6em]
  \item \textbf{Threshold Scatter} (page 1)
  \item \textbf{Value Histogram} (page 2)
  \item \textbf{Summary Report} (lot/wafer counts, file metadata)
  \item \textbf{8-Condition Distribution} bar chart
  \item \textbf{Per-lot wafer grids} — a $5\times5$ grid (\code{MAX\_WAFERS\_PER\_PAGE}$=25$) per page
\end{enumerate}

Each grid page is a self-contained single-page \emph{vector} PDF produced by a
picklable worker (\code{\_render\_grid\_page}); the main process merges them in
order, so parallelism costs no quality.

\section{PPTX}
\code{create\_powerpoint\_report} is fed the PNG snapshots that the PDF pages
emit as a side effect. It sorts them by leading page number (numerically, so a
2000-slide deck past page 999 stays ordered), builds a title slide, and adds one
image slide per page. Consequently the deck order matches the PDF, including the
scatter and histogram front pages.

\section{CSV}
\code{export\_color\_summary\_csv} writes a per-wafer colour-zone count summary.

% ===========================================================================
\chapter{The Graphical Interface}\label{ch:gui}

\section{Main window}
\code{ui/main\_window.DataProcessorUI} is a split view: a left control panel
(input file, thresholds, orientation, output folder, run/automation buttons) and
a right panel of tabs.

\begin{longtable}{@{}p{4.2cm}p{10.3cm}@{}}
\toprule
\textbf{Tab / widget} & \textbf{Purpose}\\
\midrule
\endhead
Live Preview ($5\times5$) & Wafer thumbnails stream in as they render; resets per
  lot.\\
Threshold Scatter & Interactive scatter of the loaded parameter with log-Y
  toggle.\\
Histogram & \code{HistogramWidget} — colour-zoned distribution with log toggles
  and a stats bar.\\
Wafer Orientation & Live diagram reflecting rotation / mirror choices.\\
\bottomrule
\end{longtable}

\section{EFF loading in the GUI}
Choosing a \code{.eff} file opens the \textbf{parameter picker}
(\code{eff\_param\_dialog}): Lot/Wafer/X/Y are shown as always-loaded
coordinates; the user ticks one or more measurement parameters (filter,
select-all, ``with spec limits'' helpers). On acceptance the tool enters
\emph{EFF mode}.

\section{Per-parameter preview pager}
In EFF mode a selector bar appears above the tabs:
\begin{lstlisting}[style=plain]
[<- Prev]  [ 1/N  <parameter>  v ]  [Next ->]   <param>: 1,475 pts . 25 wafers
\end{lstlisting}
Paging (or the dropdown) loads that parameter's data on an
\code{EffPreviewWorker} background thread and updates the scatter + histogram.
Results are cached by column index, so revisiting a parameter is instant. On
navigation the tool also:
\begin{itemize}[leftmargin=1.4em]
  \item auto-fills \textbf{Limit 1 / Limit 2} from that parameter's EFF spec
        limits, and
  \item shows the parameter name in the chart titles and a label in the
        Parameters panel.
\end{itemize}

\section{Run \& worker signals}
``Generate Report'' dispatches to \code{EffReportWorker} (EFF mode) or
\code{ReportWorker} (txt/csv). Both are \code{QThread}s emitting
\code{progress}, \code{wafer\_progress}, \code{status\_message},
\code{wafer\_ready}, an error signal, and a completion signal
(\code{report\_done} / \code{all\_done}). The window connects these to the two
progress bars, the live grid, and the status line, and can cancel mid-run.

% ===========================================================================
\chapter{Automation}\label{ch:auto}

\section{Saving \& replaying jobs}
``Save Job'' gathers the current inputs into a \code{WaferJob} and writes a
\code{.wtjob}. \code{automation/runner.py} replays it headlessly:
\code{run\_job} $\to$ \code{\_run\_loaded}, dispatching to \code{process\_eff}
for \code{.eff} inputs or \code{generate\_report} for txt/csv.

\section{Exit codes}
The runner returns a distinct exit code per outcome, surfaced to Windows Task
Scheduler as \code{LastTaskResult}:

\begin{longtable}{@{}rp{4.2cm}p{7.6cm}@{}}
\toprule
\textbf{Code} & \textbf{Constant} & \textbf{Meaning}\\
\midrule
\endhead
0  & \code{EXIT\_OK} & report generated\\
1  & \code{EXIT\_ERROR} & unexpected error during generation\\
2  & \code{EXIT\_USAGE} & no job path given\\
3  & \code{EXIT\_NO\_OUTPUT} & ran but produced no PDF\\
10 & \code{EXIT\_MISSING\_INPUT} & the job's input file (e.g.\ \code{.eff}) is missing\\
11 & \code{EXIT\_NO\_MATCH\_IN\_FOLDER} & folder-mode job matched no file\\
12 & \code{EXIT\_BAD\_JOB} & the \code{.wtjob} is missing / invalid\\
\bottomrule
\end{longtable}

\section{Windows Task Scheduler integration}
\code{automation/scheduler.py} registers a task named \code{WaferTool\_Job\_<name>}
per job. Daily and weekly triggers use PowerShell \code{ScheduledTasks} cmdlets;
monthly uses Task Scheduler XML (PowerShell has no monthly-trigger cmdlet). Tasks
run as the logged-on user with no stored password. The command line is either the
frozen executable (\code{App.exe -{}-run-job <job>}) or
\code{python run\_job.py <job>} in development.

\section{Schedule Jobs dialog \& Status column}
\code{scheduler\_dialog.SchedulerDialog} lists the registered tasks and lets the
user schedule / run-now / remove them. A colour-coded \textbf{Status} column maps
the last exit code (via \code{\_status\_for}) to a readable outcome:

\begin{longtable}{@{}p{6.8cm}p{2.4cm}p{4.6cm}@{}}
\toprule
\textbf{Status} & \textbf{Colour} & \textbf{Trigger}\\
\midrule
\endhead
\textcolor{zgreen}{$\checkmark$ Success} & green & code 0\\
\textcolor{zred}{$\times$ Missing input file (EFF)} & red & code 10\\
\textcolor{zamber}{$\triangle$ No matching file in folder} & amber & code 11\\
\textcolor{zred}{$\times$ Bad / missing job file} & red & code 12\\
\textcolor{zred}{$\times$ Ran but no output} & red & code 3\\
\textcolor{zred}{$\times$ Failed (error)} & red & code 1 / other\\
$\blacktriangleright$ Running\ldots & blue & Task Scheduler ``running''\\
--- Not run yet & grey & never run\\
\bottomrule
\end{longtable}

% ===========================================================================
\chapter{Concurrency Model}\label{ch:concurrency}

\begin{itemize}[leftmargin=1.4em]
  \item \textbf{Qt threads} — long operations (report generation, EFF parameter
        preview) run on \code{QThread}s so the UI never blocks. They communicate
        with the UI only through Qt signals.
  \item \textbf{Process pools} — CPU-bound rendering (wafer PNGs, PDF grid pages)
        fans out across processes with \code{ProcessPoolExecutor}; results are
        consumed in submission order so the live preview stays ordered.
  \item \textbf{Cancellation} — a \code{threading.Event} in each worker; the
        pipeline observes it through the progress callback (raising
        \code{RuntimeError}) and at chunk / parameter boundaries.
  \item \textbf{Within vs.\ across parameters} — the EFF pipeline processes
        parameters sequentially, each internally parallel, to keep all cores busy
        without oversubscribing nested pools.
\end{itemize}

% ===========================================================================
\chapter{Performance \& Scalability}\label{ch:perf}

\section{What scales well}
EFF extraction \emph{streams} (no whole-file materialisation); per-wafer grids are
freed immediately after drawing; rendering is parallelised; the PPTX slide images
are produced as a by-product of the PDF pages rather than a third pass; and
uniform wafers are saved at a lower DPI (see \S\,\ref{ch:render}) so common
single-colour wafers cost little disk.

\section{Known limits}
\begin{enumerate}[leftmargin=1.6em]
  \item \textbf{No wafer-count ceiling.} Everything downstream is
        $O(n_\text{wafers})$: the statistics loop, the CSV loop, $N$ PDF pages,
        $N$ PNGs, $N$ slides. A file that resolves to $\sim$195{,}000 wafers
        (dense parametric data across many lots) makes a run impractical. The
        recommended fix is a \emph{guardrail} that warns before rendering an
        unreasonable wafer count.
  \item \textbf{Pure-Python per-wafer loops} in \code{statistics.py} and
        \code{csv\_exporter.py} iterate every group in Python; vectorising with a
        single pandas \code{groupby().agg()} would cut these dramatically at 100k+
        wafers.
  \item \textbf{Full-file materialisation on the txt path} — the converted-text
        loader reads the entire file into memory.
  \item \textbf{Interpolation is cheap for sparse wafers.} A measured experiment
        showed that caching interpolated grids to disk to avoid a second
        interpolation was \emph{slower} than recomputing (disk load
        $\approx52$\,ms vs.\ \code{griddata} $\approx11$\,ms for $\sim$9-point
        wafers), so that optimisation was deliberately \emph{not} adopted.
\end{enumerate}

% ===========================================================================
\chapter{Modularity Assessment}\label{ch:modularity}

\textbf{Strengths.} Clear downward-only layering; the EFF layer
(\code{eff\_loader} $\to$ \code{eff\_pipeline} $\to$ \code{eff\_worker}) cleanly
separates I/O, GUI-free orchestration, and the Qt adapter; exporters are isolated
behind functions; one core renders for both GUI and report.

\textbf{Opportunities.}
\begin{enumerate}[leftmargin=1.6em]
  \item \code{ui/main\_window.py} is large ($\sim$1.6k lines) and mixes reusable
        canvases, the main window, and business logic
        (\code{\_gather\_job}, EFF wiring). Extracting the canvases to their own
        module and moving job-gathering out of the UI would slim it.
  \item The scatter/histogram exist \emph{twice} (GUI canvases vs.\
        \code{report\_charts}); sharing one axes-level renderer would remove drift.
  \item Some dead / legacy code (an unused \code{HistogramCanvas},
        \code{*\_old.py} files) can be removed.
\end{enumerate}

% ===========================================================================
\chapter{Build \& Run}\label{ch:build}

\begin{itemize}[leftmargin=1.4em]
  \item \textbf{Interactive:} \code{python main.py} (PyQt6).
  \item \textbf{Headless job:} \code{python run\_job.py <job.wtjob>} — fixes
        \code{sys.path} then calls \code{automation.runner.main}.
  \item \textbf{Packaging:} PyInstaller spec files (\code{main.spec},
        \code{WaferMapTool.spec}); the PPTX template and compiled parser are
        bundled and resolved from \code{sys.\_MEIPASS} when frozen.
  \item \textbf{Dependencies:} PyQt6, matplotlib, numpy, pandas, scipy, pypdf,
        python-pptx (optional; PPTX is skipped gracefully if absent).
\end{itemize}

\vspace{1cm}
\begin{center}
\textcolor{muted}{\rule{0.5\textwidth}{0.4pt}}\\[4pt]
\textcolor{muted}{\small End of document}
\end{center}

\end{document}
